\begin{claim}{D.1.7}{Proved}{The simply connected odd-primary Moore targets}
Let $p$ be odd and let $n\geq2$. For the usual Moore space
$M(\ZZ/p,n)=S^n\cup_p e^{n+1}$, its single-target null class is
closed under desuspension. There is no finiteness hypothesis on the
source. More precisely, for any $d\geq3$ and any pointed space $X$,
\[
 \Map_*(X,\Omega P^d(p))\simeq *
 \quad\Longrightarrow\quad
 \Map_*(X,P^e(p))\simeq *\quad\hbox{for every }e\geq3,
 \tag{D.1.7.1}
\]
where $P^e(p)=S^{e-1}\cup_p e^e$ is indexed by its top cell.
\end{claim}
\begin{proof}
Fix $X$ and write
\[
 \mathcal R_X=\{Y:\Map_*(X,Y)\simeq *\}.
\]
This class is closed under homotopy limits, loops and homotopy
retracts. The proof has three steps: obtain all sufficiently high
looped Moore targets; use a convergent loop--suspension resolution
to remove their loops; and use finite cone cubes to lower their
dimensions. All the closure statements concern the whole based
mapping space, including its set of components.

\medskip\noindent
\textbf{The precise retraction input and the bottom-dimensional case.}
For odd $p$ the smash splitting is
\[
 P^a(p)\wedge P^b(p)
       \simeq P^{a+b}(p)\vee P^{a+b-1}(p),\qquad a,b\geq2.
 \tag{D.1.7.2}
\]
It is the splitting used in the proof of
\cite[Lemma~1.5, pp.~13--14]{I_CMN1987}; the hypotheses and formula
are also explicitly recorded in \cite[Lemma~4.5]{I_Boyde2021}.
We use two retraction theorems from the original CMN paper:
\[
 \Omega P^{2a+2}(p)
 \simeq S^{2a+1}\{p\}\times
       \Omega\bigvee_{j\geq0}P^{4a+2ja+3}(p),
       \qquad a\geq1,                              \tag{D.1.7.3}
\]
and the assertion that a mod-$p$ homology monomorphism between
loop spaces of bouquets of simply connected mod-$p$ Moore spaces
has a left homotopy inverse. These are respectively
\cite[Proposition~1.1, p.~9, and Proposition~1.7, p.~15]{I_CMN1987}.
The symbol $S^{2a+1}\{p\}$ denotes the fibre of the degree-$p$
map on the $p$-local sphere; only the Moore-bouquet factor will be
used. CMN's retractions are $p$-local statements. Every simply
connected finite mod-$p$ Moore space, and every finite wedge or
smash of these spaces, has finite $p$-primary homotopy groups by
the Serre-class theorem. These spaces are therefore already
$p$-local. The retractions used here consequently apply to their
ordinary homotopy types. The displayed range $a\geq1$ was checked
on the printed page of Proposition~1.1, including its endpoint.

For clarity we do not use CMN's odd-dimensional Corollary~1.9
outside its printed range $a\geq2$. In particular the case
$P^3(p)$ is established as follows. Let $d\geq3$ be odd, set
$A=P^{d-1}(p)$, and let
\[
 [\id,\id]:\Sigma(A\wedge A)\longrightarrow\Sigma A=P^d(p)
\]
be the generalized Whitehead product. The lower Moore summand in
(D.1.7.2), before suspension, gives a map
$j:P^{2d-3}(p)\to A\wedge A$. Put
\[
 h=[\id,\id]\circ\Sigma j:P^{2d-2}(p)\longrightarrow P^d(p).
\]
We check that $\Omega h$ is a mod-$p$ homology monomorphism.
Write
\[
 H_*(\Omega P^d(p);\FF_p)=T(u,v),\qquad
 |u|=d-2,\quad |v|=d-1,\quad\beta v=u.
\]
This tensor-algebra formula is the James theorem for the connected
space $A$. The adjoint Whitehead product has the graded commutator
as its homology map. Since $|u|$ is odd and $|v|$ is even, the
two generators of $H_*(\Omega P^{2d-2}(p);\FF_p)$ have images,
after choosing the compatible Moore orientations,
\[
                  2u^2,\qquad vu-uv.                \tag{D.1.7.4}
\]
Here is a check that the second formula is independent of the
chosen complementary summand. Normalize the lower generator of
$j$ to $u\otimes u$. Its upper generator then maps to
$c(v\otimes u)+b(u\otimes v)$ with $c-b=1$, since
$\beta(v\otimes u)=u\otimes u$ and
$\beta(u\otimes v)=-u\otimes u$. The commutator sends this
combination to $(c-b)(vu-uv)$. It sends $u\otimes u$ to $2u^2$.
Thus the two images also satisfy the required Bockstein relation.

These images freely generate a tensor subalgebra. Indeed, with
the lexicographic order $v>u$, their leading words are $uu$ and
$vu$, respectively, with nonzero coefficients. Both words have
length two. A concatenation of them has a unique division into
two-letter blocks, so distinct monomials in the two proposed
generators have distinct leading words. This proves injectivity.
CMN Proposition~1.7 therefore makes $\Omega P^{2d-2}(p)$ a
homotopy retract of $\Omega P^d(p)$. Its source and target are
loop spaces of simply connected Moore spaces even when $d=3$;
the proposition has no additional connectivity restriction here.
Consequently an even-dimensional looped Moore space belongs to
$\mathcal R_X$ whenever the original $\Omega P^d(p)$ does.

\medskip\noindent
\textbf{All sufficiently high looped Moore spaces.}
Apply (D.1.7.3) to that even-dimensional Moore space, and retain
two bouquet summands. For some integers $a_1,a_2\geq3$ we obtain
\[
 \Omega(P^{a_1}(p)\vee P^{a_2}(p))\in\mathcal R_X.
 \tag{D.1.7.5}
\]
For example the first two summands of (D.1.7.3) have top
dimensions $4a+3$ and $6a+3$.
The Hilton--Milnor factor indexed by the basic commutator with
$k$ occurrences of the first letter and one of the second is
\[
 \Omega\Sigma\bigl(
      P^{a_1-1}(p)^{\wedge k}\wedge P^{a_2-1}(p)\bigr),
      \qquad k\geq1.
\]
It is a homotopy retract of (D.1.7.5). Repeated use of
(D.1.7.2) makes every $\Omega P^D(p)$ with
\[
        k(a_1-2)+a_2\ \leq D\leq\ k(a_1-1)+a_2
 \tag{D.1.7.6}
\]
a retract of this factor. Each integer in the interval occurs:
the smash of $k+1$ Moore spaces has a Moore summand with each of
the $k+1$ possible dimension reductions $0,1,\ldots,k$.
The lower endpoint for $k+1$ minus the upper endpoint for $k$
is $a_1-2-k$. Thus the integer intervals overlap or are adjacent
for all sufficiently large $k$, and their union contains every
integer above some bound $N$. We have proved
\[
                  \Omega P^D(p)\in\mathcal R_X
                  \quad(D\geq N).                  \tag{D.1.7.7}
\]
The Hilton--Milnor theorem and its Moore-smash application used
here are also given in \cite[Proposition~2.2 and its proof,
pp.~18--19]{I_CMN1987}; the interval calculation is the extra
step which gives every sufficiently large dimension.

We will need a slightly stronger consequence. If $W$ is a finite
wedge of Moore spaces of top dimension at least $N$, every
Hilton--Milnor factor of $\Omega W$ splits into looped Moore
spaces of top dimension at least $N$. A word of length $q$ has
minimum possible top dimension
$\sum_{i=1}^q e_i-2q+2\geq q(N-2)+2\geq N$.
The factors' connectivities increase with word length, so the
weak product supplied by Hilton--Milnor is weakly equivalent to
the ordinary product. Products and retracts are allowed in
$\mathcal R_X$, proving $\Omega W\in\mathcal R_X$.
The same conclusion holds for a finite-type wedge: truncate it
by top dimension, use the projections that collapse the tails,
and take the homotopy limit. In each fixed homotopy degree this
tower is eventually constant, so its homotopy limit is the
original wedge. This argument uses target connectivity, not
compactness of $X$.

\medskip\noindent
\textbf{Removing loops and lowering dimensions.}
Let $K$ be any finite wedge of Moore spaces of top dimension at
least $N$, and put $T=\Omega\Sigma$. The iterated James splitting
from C.9 gives
\[
 \Sigma T^rK\simeq
       \bigvee_{m\geq1}\bigvee^{c_{r,m}}\Sigma K^{\wedge m},
       \qquad r\geq0,                              \tag{D.1.7.8}
\]
where each multiplicity $c_{r,m}$ is finite. To recall why,
$\Sigma TV\simeq\bigvee_{q\geq1}\Sigma V^{\wedge q}$;
distributing smash products and moving the single suspension
between factors gives the induction. At a fixed weight $m$
there are only finitely many ordered decompositions of $m$ into
positive weights, so all the multiplicities remain finite.
By (D.1.7.2), the right side of (D.1.7.8) is a finite-type wedge
of Moore spaces of top dimension at least $N+1$. Its loop space
therefore belongs to $\mathcal R_X$ by (D.1.7.7) and the preceding
paragraph. Thus $T^{r+1}K\in\mathcal R_X$ for every $r\geq0$.

For every simply connected pointed space the derived
loop--suspension resolution converges:
\[
             K\simeq\holim_{[q]\in\Delta}T^{q+1}K.
 \tag{D.1.7.9}
\]
This is exactly \cite[Theorem~1.8(b)]{FGHBlomquistHarper2017},
with $r=1$ in that paper, or
\cite[Sections~7.9 and~9.9]{FGHBousfield2003}. There is no
finite-type hypothesis on the mapping source. Closure under
homotopy limits proves that $\mathcal R_X$ contains every such
finite wedge $K$.

Finally we spell out the desuspension mechanism, so that this
step does not assume the property being proved. If $V$ is simply
connected and every finite wedge of $\Sigma V$ lies in a
homotopy-limit-closed class $\mathcal R$, form the strongly
homotopy cocartesian $s$-cube generated by $s$ copies of
$V\to CV$. Each noninitial vertex is a finite wedge of
$\Sigma V$, including the empty wedge. Let $Y_s$ be the homotopy
limit of the punctured cube. Then $Y_s\in\mathcal R$.
The higher Blakers--Massey theorem makes the comparison
$V\to Y_s$ increasingly connected as $s$ increases. The cone
cubes have compatible projection maps, and hence
$V\simeq\holim_sY_s\in\mathcal R$. This is the Barratt
argument in \cite[Theorem~8, p.~7]{I_Strom2012}.

Apply this argument first to every finite wedge of $P^{N-1}(p)$,
whose suspension is a finite wedge of $P^N(p)$. It puts all
these wedges in $\mathcal R_X$. Repeat, one dimension at a time,
down to top dimension three. Together with the already proved
high-dimensional assertion, this proves (D.1.7.1). Taking
$e=d=n+1$ and using the suspension--loop adjunction proves the
claim.
\end{proof}

The odd-prime hypothesis is used both in (D.1.7.2) and in the
nonzero coefficient $2$ in (D.1.7.4). Neither step gives the
mod-$2$ Moore assertion by substitution.

\begin{claim}{D.1.8}{Proved}{The simply connected mod-two Moore targets}
For every $n\geq2$, the single-target null class of
$M(\ZZ/2,n)$ is closed under desuspension, with no restriction on
the source space. More strongly, nullity into any
$\Omega P^d(2)$, $d\geq3$, implies nullity into every
$P^e(2)$, $e\geq3$. The additional geometric input is the explicit
homotopy retraction
\[
 \Omega\bigl(P^{11d-15}(2)\vee P^{11d-15}(2)\bigr)
       \quad\hbox{of}\quad \Omega P^d(2),\qquad d\geq4.
 \tag{D.1.8.1}
\]
For the bottom case one can first use the retraction of
$\Omega P^6(2)$ from $\Omega P^3(2)$.
\end{claim}
\begin{proof}
We give the two-summand construction, including a reproducible
finite certificate for its only larger calculation. It does not
use an odd-primary splitting of Moore smash products.

\medskip\noindent
\textbf{The geometric maps.}
Fix $d\geq4$ and write
\[
 A=P^{d-1}(2),\qquad |u|=d-2,\quad |v|=d-1,\qquad
 s=|u|+|v|=2d-3,
\]
where $u,v$ generate $\widetilde H_*(A;\FF_2)$.
The space $A$ is a simply connected finite suspension and is
already $2$-local. For odd $m$, let
$\widetilde L_m(A)$ be the telescope of the
Dynkin--Specht--Wever map $\frac1m\beta_m$ on
$\Sigma A^{\wedge m}$. There are maps
\[
 p_m:\Sigma A^{\wedge m}\longrightarrow\widetilde L_m(A),
 \qquad
 i_m:\widetilde L_m(A)\longrightarrow\Sigma A^{\wedge m}
\]
such that $p_mi_m\simeq\id$ and, on mod-$2$ homology,
\[
 (i_mp_m)_*(x_1\cdots x_m)
       =\beta_m(x_1\cdots x_m)
       =[[\cdots[x_1,x_2],x_3],\ldots,x_m].
 \tag{D.1.8.2}
\]
Suspension symbols on homology are suppressed here and below.
This splitting is constructed by the two complementary
telescopes and suspension comultiplication in
\cite[Section~2, pp.~3248--3249]{I_ChenWu2013}.
It does not require the representing map to be a co-H map.
Since $A$ is a suspension, \cite[Theorem~1.6]{I_Wu1998}
makes $\Omega\widetilde L_{11}(A)$ a homotopy retract of
$\Omega\Sigma A$. The sequence in that theorem can be the odd
primes; its conditions exclude divisibility by $2$ and divisibility
of one chosen index by another. Thus the index $11$ is permitted.

There is a map
\[
       \alpha:S^s\longrightarrow A\wedge A,
       \qquad \alpha_*(\iota_s)=uv+vu,
 \tag{D.1.8.3}
\]
by \cite[Proposition~2.1, pp.~3249--3250]{I_ChenWu2013}.
The map
$e_3=p_3\Sigma(\id_A\wedge\alpha):
\Sigma^{s+1}A\to\widetilde L_3(A)$
is a homotopy equivalence. To verify it in this case, the image
of $\beta_3$ on $\widetilde H_*(\Sigma A^{\wedge3};\FF_2)$
has the two basis elements $uuv+vuu$ and $uvv+vvu$;
$e_3$ takes the two generators to these elements. Thus it is a
mod-$2$ homology equivalence between simply connected spaces
whose integral homology is finite $2$-primary torsion in bounded
degrees. It is an integral homology equivalence and hence a
homotopy equivalence. Equivalently this is the first step of
\cite[proof of Proposition~3.1, p.~3251]{I_ChenWu2013}.
Define $r_3=e_3^{-1}p_3$. Its homology map has the complete table
\[
\begin{array}{c|cccccccc}
 w&uuu&uuv&uvu&uvv&vuu&vuv&vvu&vvv\\ \hline
 (r_3)_*(w)&0&0&u&v&u&v&0&0.
\end{array}                                      \tag{D.1.8.4}
\]

Iterate this actual projection, moving suspension factors through
smash products, to obtain
\[
 \varphi_k:\Sigma A^{\wedge(2k+1)}
                     \longrightarrow\Sigma^{ks+1}A,
 \qquad
 \varphi_1=r_3,\qquad
 \varphi_k=(\Sigma^{(k-1)s}r_3)
                (\varphi_{k-1}\wedge\id_A\wedge\id_A).
 \tag{D.1.8.5}
\]
Its homology map applies the table (D.1.8.4) successively to the
first three letters, to the result and the next two letters, and
so on. This also describes what happens when an intermediate
value is zero: the final value is zero.

Set $B=\Sigma^{5s+1}A=P^{11d-15}(2)$ and
$L=\widetilde L_{11}(A)$. Let $\sigma=(1\ 8)$ and
$\tau=(1\ 3)$ act on the eleven smash factors. With
$j_0=\Sigma(\id_A\wedge\alpha^{\wedge5}):
B\to\Sigma A^{\wedge11}$, define
\[
 j_1=p_{11}j_0,\quad j_2=p_{11}\sigma j_0,
 \qquad
 r_1=\varphi_5i_{11},\quad r_2=\varphi_5\tau i_{11}.
 \tag{D.1.8.6}
\]
All four maps have now been specified geometrically.

\medskip\noindent
\textbf{The exact finite calculation.}
For either generator of $\widetilde H_*(B;\FF_2)$ the matrix
of the four composites $r_i j_j$ is
\[
                     \begin{pmatrix}1&1\\1&0\end{pmatrix}.
 \tag{D.1.8.7}
\]
Here are explicit finite instructions and intermediate integer
counts that verify the matrix. They also fix all permutation and
commutator conventions, so that no representation-theoretic
decomposition is being assumed. Write $u=0$, $v=1$ and put
\[
 w(\epsilon,x)=
 (x,\epsilon_1,1-\epsilon_1,\ldots,
       \epsilon_5,1-\epsilon_5),
 \qquad\epsilon\in\{0,1\}^5,\quad x\in\{0,1\}.
\]
For a word $z=(z_1,\ldots,z_{11})$ and
$S\subseteq\{2,\ldots,11\}$, let $D_S(z)$ be the word obtained
by listing the letters with indices in $S$ in decreasing order,
then $z_1$, then the remaining letters in increasing order.
Expansion of the successive commutators gives the exact identity
\[
               \beta_{11}(z)=\sum_{S\subseteq\{2,\ldots,11\}}
                                  D_S(z)\quad\hbox{over }\FF_2.
 \tag{D.1.8.8}
\]
Indeed at each commutator step the new letter is put either on
the left or on the right; the former choices accumulate in
reverse order and the latter in their original order.

Let $f(z)$ be the repeated ternary reduction in (D.1.8.4), with
value $\bot$ if a zero is encountered. For
$(\rho_1,\rho_2)=(\id,\tau)$ and
$(\gamma_1,\gamma_2)=(\id,\sigma)$ define the integer counts
\[
 C_{ij}(x)=
 \sum_{\epsilon\in\{0,1\}^5}
 \sum_{S\subseteq\{2,\ldots,11\}}
  \mathbf1\{f(\rho_iD_S(\gamma_jw(\epsilon,x)))=x\}.
\]
Applying the eight-entry table five times gives
\[
             C(0)=C(1)=
                  \begin{pmatrix}3125&3125\\3039&2460\end{pmatrix}.
 \tag{D.1.8.9}
\]
Reduction modulo two is (D.1.8.7). Each entry has exactly
$2^5\cdot2^{10}$ candidate terms, so this is a bounded exhaustive
calculation. A second calculation cancels equal words after
each commutator in (D.1.8.2) before doing any projections; the
two embedding supports then have sizes $62$ and $126$, and
give the same matrix. The complete short verifier
\texttt{code/I\_lie\_moore.py} implements both formulas
independently, compares them on every input word, and checks
the displayed integer and mod-$2$ matrices. Its exact output
is \texttt{code/outputs/I\_lie\_moore.json}.

\medskip\noindent
\textbf{From the matrix to the wedge retraction.}
Being a retract of a suspension, $L$ has a comultiplication
\[
 \Delta_L=(p_{11}\vee p_{11})
                \Delta_{\Sigma A^{\wedge11}}i_{11}:L\to L\vee L
\]
whose two projections are homotopic to the identity. Thus the
two projections in (D.1.8.6) give one actual map to a wedge:
\[
 R=(r_1\vee r_2)\Delta_L:L\longrightarrow B\vee B,
 \qquad J=j_1\vee j_2:B\vee B\longrightarrow L.
\]
By (D.1.8.7), $RJ$ is a mod-$2$ homology equivalence. It is a
homotopy equivalence: its mapping cone is simply connected,
has finite $2$-primary integral homology, and has zero mod-$2$
homology, hence zero integral homology. Replacing $R$ by
$(RJ)^{-1}R$ proves that $B\vee B$ is a homotopy retract of
$L$. Looping, and using the Wu retraction above, proves
(D.1.8.1).

For $d=3$ use the explicit Moore-space retraction
\[
                  \Omega P^6(2)\quad\hbox{of}\quad\Omega P^3(2).
 \tag{D.1.8.10}
\]
It is the case $n=2$, $k=1$ of the formula
$\Omega P^{n+1+k(2n-1)}(2)\to\Omega P^{n+1}(2)$ stated
in \cite[proof of Corollary~1.4, p.~3258]{I_ChenWu2013},
where the printed range is $n\geq2$, $k\geq1$.
We only need this endpoint, not the broader prime-indexed product
formula in that paper's Theorem~1.3. Applying (D.1.8.1) with
$d=6$ then supplies the required two equal summands. This use of
the paper's explicit $P^3(2)$ application does not identify an
odd-degree telescope of $\RP^2$ with a CW complex having finitely
many ordinary cells in each degree; that identification would
be false, since the telescope has $\pi_2=\ZZ_{(2)}$.

\medskip\noindent
\textbf{Why two equal summands suffice.}
Suppose now $\Omega P^d(2)\in\mathcal R_X$, where $d\geq3$.
We have just obtained an integer $D>d$ with
\[
       \Omega(P^D(2)\vee P^D(2))\in\mathcal R_X.
\]
Put $K=P^{D-1}(2)$. The Hilton--Milnor word containing $q-1$
copies of the first letter and one copy of the second makes
$\Omega\Sigma K^{\wedge q}$ a retract for every $q\geq2$;
the case $q=1$ is either original wedge summand. Consequently
all these loop spaces belong to $\mathcal R_X$.

If $V$ is any finite wedge of copies of $K$, the same theorem
shows that the loops on every finite wedge of spaces
$\Sigma V^{\wedge m}$ belong to $\mathcal R_X$: all its factors
have the form $\Omega\Sigma K^{\wedge q}$. Factor connectivities
increase with word length, so the weak product is weakly
equivalent to the product. The iterated James splitting
(D.1.7.8), now used without any Moore smash splitting, expresses
$\Sigma(\Omega\Sigma)^rV$ as a finite-type wedge of precisely
these smash powers. Its truncation tower has increasingly
connected tails. Therefore all
$(\Omega\Sigma)^{r+1}V$ belong to $\mathcal R_X$.
The convergent simply connected resolution (D.1.7.9) proves
$V\in\mathcal R_X$.

We have proved nullity for every finite wedge of $P^{D-1}(2)$.
The finite cone-cube argument at the end of D.1.7 lowers this
whole family by one dimension at a time, down to top dimension
three. In particular $P^d(2)\in\mathcal R_X$, as $3\leq d<D$.
Moreover $P^{D-1}(2)$ and its loop space now belong to
$\mathcal R_X$. Repeating the same construction with $d=D-1$
increases the upper bound to $11(D-1)-15$, and iteration makes
these bounds unbounded. Thus every $P^e(2)$ with $e\geq3$
belongs to $\mathcal R_X$, proving the stronger assertion.
Adjunction proves the requested desuspension closure. The only
homotopy limits taken in this proof are limits in the targets;
the source $X$ has been arbitrary throughout.
\end{proof}

\begin{claim}{D.1.9}{Proved}{The Moore spaces with cyclic fundamental group}
For odd $p$, the single-target null class of
$M(\ZZ/p,1)=P^2(p)$ is closed under desuspension, for arbitrary
source spaces. For $p=2$ this assertion is false:
$P^2(2)=\RP^2$ admits an explicit simply connected
factorial-telescope counterexample. Consequently, among the
usual Moore targets $M(\ZZ/p,n)$ with $p$ prime and $n\geq1$,
the only failure is $M(\ZZ/2,1)$.
\end{claim}
\begin{proof}
Let $p$ be odd, put $Y=P^2(p)$, and suppose
$\Map_*(X,\Omega Y)\simeq *$. We may replace $X$ by a CW
model. This hypothesis forces $X$ to be connected: otherwise
$S^0$ is a pointed retract of $X$, and $\Omega Y$ would be a
retract of the contractible mapping space. That is impossible,
since $\pi_0\Omega Y=\ZZ/p$.

The universal cover $V$ of $Y$ satisfies
\[
                         V\simeq\bigvee^{p-1}S^2.
 \tag{D.1.9.1}
\]
Indeed $Y$ is the presentation complex of $\langle a\mid a^p\rangle$,
and its $p$-sheeted universal cover is a simply connected finite
two-dimensional complex of Euler characteristic $p$. Its only
nonzero reduced integral homology is free of rank $p-1$ in degree
two. Representatives of a homology basis under the Hurewicz
isomorphism give a homology equivalence from the displayed wedge;
the simply connected Whitehead theorem makes it a homotopy
equivalence.

Since $X$ is connected, its based maps to $\Omega Y$ land in
the base component $\Omega V$. Thus $\Omega V$ belongs to
$\mathcal R_X=\{Z:\Map_*(X,Z)\simeq *\}$.
There are at least two sphere summands in (D.1.9.1).
Hilton--Milnor on these two summands gives a retract
$\Omega S^{q+1}$ for every $q\geq1$: use a word of length
$q$ with $q-1$ occurrences of one letter and one of the other
when $q\geq2$, and a sphere summand when $q=1$.
Hence
\[
                      \Omega S^d\in\mathcal R_X
                      \qquad(d\geq2).              \tag{D.1.9.2}
\]

Here is the unlooping and finite-target step with its source
hypotheses made explicit. If $W$ is a finite wedge of spheres
of dimension at least two, iterated James splitting expresses
$\Sigma(\Omega\Sigma)^rW$ as a finite-type wedge of spheres.
Hilton--Milnor and (D.1.9.2) put the loops on every finite
truncation in $\mathcal R_X$. Increasing connectivity of the
tails and closure under homotopy limits put
$(\Omega\Sigma)^{r+1}W$ in $\mathcal R_X$ for every $r\geq0$.
The simply connected resolution (D.1.7.9) now gives
$W\in\mathcal R_X$. In particular every $S^d$, $d\geq2$,
belongs to this class.

The class $\mathcal R_X$ is a resolving kernel and is closed
under extensions by fibrations. We can therefore apply
\cite[Corollary~12(a), pp.~8--9]{I_Strom2012}, which puts every
simply connected finite complex in $\mathcal R_X$ as soon as
the sphere dimensions include the indicated infinite arithmetic
progression. Its proof uses the collection-level wedge theorem,
the cone-length theorem, and finite cone cubes. It has no
finite-type hypothesis on $X$; the finite-type source hypothesis
in that paper's Theorem~1 enters only for the stronger assertion
about arbitrary finite-dimensional targets. We only need the
two finite targets
\[
                           V,\ \Sigma Y\in\mathcal R_X.
 \tag{D.1.9.3}
\]
In particular $\Omega\Sigma Y\in\mathcal R_X$.

Consider any based map $f:X\to Y$ and the unit
$\eta_Y:Y\to\Omega\Sigma Y$. The composite $\eta_Yf$ is
nullhomotopic by (D.1.9.3). On fundamental groups the unit is
the abelianization map
\[
 \pi_1Y\longrightarrow H_1(Y;\ZZ)
     \cong H_2(\Sigma Y;\ZZ)
     \cong\pi_2\Sigma Y=\pi_1\Omega\Sigma Y.
\]
For $Y=P^2(p)$ it is an isomorphism $\ZZ/p\to\ZZ/p$.
Thus $\pi_1f=0$, so $f$ lifts to the universal cover $V$.
The lift is nullhomotopic by (D.1.9.3); hence so is $f$.
We have proved that $\Map_*(X,Y)$ has only its constant
component. Its loop space is the given contractible space
$\Map_*(X,\Omega Y)$, so all its positive homotopy groups
also vanish. Therefore $Y\in\mathcal R_X$, as required.
Notice that no loop--suspension convergence statement has been
applied to the nonsimply connected space $Y$ itself.

For $p=2$, put
\[
 X=\operatorname{Tel}\bigl(S^2\xrightarrow{2}S^2
              \xrightarrow{3}S^2\xrightarrow{4}\cdots\bigr).
\]
This is simply connected. The integral $S^2$ calculation in
D.1.2 gives
\[
 \Map_*(X,\Omega S^2)\simeq *,\qquad
 \pi_0\Map_*(X,S^2)
       =\lim{}^1\bigl(\ZZ\xleftarrow{2}\ZZ
                      \xleftarrow{3}\ZZ\xleftarrow{4}\cdots\bigr)
       =\widehat\ZZ/\ZZ\ne0.
\]
Every based map from $X$, or from a suspension of $X$, to
$\RP^2$ lifts uniquely through its universal cover $S^2$.
The corresponding based mapping-space comparison is a weak
equivalence. Thus
$\Map_*(X,\Omega\RP^2)\simeq *$ whereas
$\Map_*(X,\RP^2)\not\simeq *$. This proves the failure and,
together with D.1.7--D.1.8, the complete stated classification.
\end{proof}
